% Appendix D: Code and Software Interface
\chapter{Code and Software Interface}
\label{app:code}

\section{Software Repository Structure}

The complete software package is organized as follows:

\begin{verbatim}
uft-clifford-torsion/
|-- src/
|   |-- spectral_dimension/     # Spectral dimension evolution
|   |-- waveform_generator/     # LISA waveform templates
|   |-- black_hole/            # Black hole evolution
|   \-- cosmology/             # Cosmological simulations
|-- data/
|   |-- lisa_sensitivity/      # LISA noise curves
|   |-- gw_waveforms/          # Gravitational wave templates
|   \-- numerical_results/     # Simulation outputs
|-- tests/
|   \-- unit_tests/            # Validation tests
\-- docs/
    \-- api_documentation/     # Code documentation
\end{verbatim}

\section{Python API}

\subsection{Spectral Dimension Evolution}

\begin{verbatim}
from uft.spectral_dimension import SpectralDimension

# Initialize with parameters
sd = SpectralDimension(
    tau_0=1e-6,
    T_GUT=1e16,  # GeV
    alpha=0.1
)

# Calculate d_s at given temperature
d_s = sd.calculate(T=1e10)  # Returns dimension

# Full evolution
times, d_s_values = sd.evolve(
    T_start=1e19,  # Planck scale
    T_end=1e-4,    # BBN scale
    n_points=1000
)
\end{verbatim}

\subsection{Gravitational Waveform Generation}

\begin{verbatim}
from uft.waveform import WaveformGenerator

# Generate 6-polarization waveform
gen = WaveformGenerator(
    m1=1e5,  # solar masses
    m2=1e5,
    distance=1e3,  # Mpc
    tau_0=1e-6
)

# Get waveform in frequency domain
f, h_plus, h_cross, h_x, h_y, h_b, h_l = \
    gen.generate_frequency_domain(
        f_min=1e-4,
        f_max=1.0
    )
\end{verbatim}

\subsection{Black Hole Evolution}

\begin{verbatim}
from uft.black_hole import BlackHoleEvolution

# Initialize black hole
bh = BlackHoleEvolution(
    M_initial=10,  # solar masses
    tau_0=1e-6
)

# Evolve until remnant formation
times, masses, torsions = bh.evolve(
    t_max=1e67,  # years
    dt=1e60
)

# Check if remnant formed
if bh.has_remnant():
    M_rem = bh.remnant_mass()
\end{verbatim}

\section{Installation}

\begin{verbatim}
# Clone repository
git clone https://github.com/dpsnet/uft-clifford-torsion.git
cd uft-clifford-torsion

# Install dependencies
pip install -r requirements.txt

# Install package
pip install -e .

# Run tests
pytest tests/
\end{verbatim}

\section{Citation}

If using this code in your research, please cite:

\begin{verbatim}
@software{UFT_Code_2024,
  author = {Research Team},
  title = {CTUFT-Clifford-Torsion: Numerical codes for unified field theory},
  url = {https://github.com/dpsnet/uft-clifford-torsion},
  year = {2024}
}
\end{verbatim}
